\subsection{A direct replacement for the missing C.17 example}
\begin{theorem}[Invariant paths with a nonlocal actual union]
\label{F:C17-fence}
Let $C$ have objects $0,1$, identities, and two distinct arrows
$f,g:0\to1$, and mark all its arrows. For every small category $Y$,
the constant-endpoint localization functor is naturally
\[
 H_{0,0}(Y)=
 \Hom_{C^Y[\operatorname{Mor}(C^Y)^{-1}]}(0_Y,0_Y)
 \cong\{v:\pi_0Y\to\mathbb Z\mid v\text{ has finite image}\}.
\]
It is invariant under every nerve weak equivalence. Nevertheless,
the actual symmetric alternating path union is not local and is not
categorically homotopy equivalent to a local category. Its specified
comparison on an infinite discrete test omits functions of infinite
image. This supplies the two C.17 assertions used in C.28 directly.
\end{theorem}
\begin{proof}
\textbf{The cover and the exact all-small-test computation.}
Let $P$ be the integer fence: its only nonidentity arrows are
$2r\to2r-1$ and $2r\to2r+1$. Send an even vertex to $0$, an odd
vertex to $1$, the right-going arrow to $f$ and the left-going arrow
to $g$. This functor $p:P\to C$ has unique lifts of arrows with
either a prescribed source or a prescribed target. The same is
true for nerve simplices, so $Bp$ is a covering map. Its deck
transformations are translations by even integers; $BP$ is the
real line with this cellular subdivision. In particular $BC$ is
the two-edge circle, not merely a space with the same homology.

For an arbitrary, possibly disconnected or empty, $Y$, let
$\operatorname{Fun}_b(Y,P)$ be the full subcategory of functors
whose integer-valued object images are uniformly bounded. Every
finite nerve diagram in it lies in
$\operatorname{Fun}(Y,P_{[-N,N]})$ for some $N$. The finite fence
$P_{[-N,N]}$ contracts categorically to $0$, fixing $0$: delete an
endpoint other than $0$, retract it onto its unique neighbor, and
repeat. At each deletion the retraction is joined to the identity
by a natural transformation in one of the two directions.
Postcomposition gives the same contraction of the functor category,
for this arbitrary $Y$. Consequently
$B\operatorname{Fun}_b(Y,P)$ is weakly contractible. Indeed its
filtered union is the union of these subcategories, and finite
simplicial sphere and homotopy data occur at a common finite stage.

A natural transformation of diagrams in $C$ lifts uniquely after a
lift of either endpoint diagram is chosen. Naturality follows by
uniqueness of the lifted composites in each square. Every component
of such a lift changes its integer height by at most one. It
therefore preserves uniform boundedness. A finite zigzag starting
at the constant $0$ diagram has a bounded lift. Conversely a bounded
lift contracts in a finite fence, so its image is in that same
component of $\operatorname{Fun}(Y,C)$. We have thus obtained a
categorical covering onto exactly the component of $0_Y$.

Two bounded lifts differ by a translation $2v(y)$, constant on
every connected component of $Y$; the function $v$ is uniformly
bounded. Conversely each such function acts freely on the cover,
and these transformations act transitively on every fiber. Its
deck group is therefore the bounded integer functions on $\pi_0Y$.
Bounded and finite-image have the same meaning for a subset of
$\mathbb Z$. The weakly contractible cover computes the fundamental
group at $0_Y$ as this group. Inverting all arrows in a category
gives its fundamental groupoid, which proves the displayed formula.
The identification sends a loop of diagrams to its componentwise
winding number and hence commutes with every restriction in $Y$.
For empty $Y$ both sides are the one-element group. Since weak
equivalences preserve $\pi_0$, this proves global invariance.

\textbf{The actual full-inclusion union.}
Let $I_n$ have vertices $-2n,\ldots,2n$ and its fence orientations,
and let $K_n$ be its endpoint-fixed path category in $C$, with both
endpoints $0$ and arbitrary interior transformations. The type is
$(+-)^{2n}$. Extend paths and transformations by the constant $0$
on both new ends to obtain full inclusions $K_n\subset K_{n+1}$.
The types are universally cofinal: their alternating run counts
are unbounded, or one may allocate each edge of any fixed type to
a matching edge in a sufficiently long alternating word.
Write $U=\bigcup_n K_n$. Lifting a path from its left endpoint
$0\in P$ gives a right endpoint $2w$. This integer $w$ is unchanged
by every endpoint-fixed transformation. The component computation
in Proposition~\ref{F:C20-paths}, applied to the point, identifies
$\pi_0U$ with $\mathbb Z$ by this same winding number.

The winding-zero full subcategory $U_0$ has a particularly literal
description. Its objects are order-preserving maps $h:P\to P$
equal to $0$ outside a finite interval. Its arrows are pointwise
order relations. For each $N$, let $Q_N$ be the full subcategory
of such maps with values in $[-N,N]$. The preceding finite-fence
contraction fixes $0$, so postcomposition contracts $Q_N$ and
preserves the finite-support condition. Thus
$U_0=\bigcup_NQ_N$ is weakly contractible and is connected.

\textbf{The obstruction to a finite categorical homotopy.}
Suppose a zigzag of $l$ natural transformations joined $1_{U_0}$
to a constant functor with value $c$. Evaluate this zigzag at an
object $h$ and then at the fence coordinate $0$. Each arrow in
$P$ changes integer height by at most one, so
\[
 |h(0)-c(0)|\le l
 \qquad\text{for every object }h\text{ of }U_0.
\]
For every positive integer $r$, however, the map
\[
 h_r(j)=\max\{0,\,2r-|j|\}
\]
is such an object and has $h_r(0)=2r$. On its support the parity
of $h_r(j)$ is the parity of $j$, and adjacent values differ by
one; across the support boundary the same fence relation or an
identity applies. These facts check every functor equation and
give the contradiction as $r$ increases.

If $U$ were local, the inclusion $U_0\to U$ and a constant functor
in its component would be joined by a finite zigzag, because
$BU_0$ is contractible. Every functor and transformation in that
zigzag remains in the same component $U_0$, contrary to the last
paragraph. Locality is preserved by categorical homotopy
equivalences, so $U$ cannot be equivalent in that sense to a local
category either. This is the asserted obstruction for the actual
symmetric union, in addition to the separate union counterexample
of Proposition~\ref{F:C28-union-counterexample}.

\textbf{The precise failed comparison.}
For the discrete category $Y=\mathbb N$, the space $BY$ is discrete
and $[BY,BU]=(\pi_0U)^{\mathbb N}=\mathbb Z^{\mathbb N}$.
Evaluation at each point and the winding computation show that
the specified map $H_{0,0}(Y)\to[BY,BU]$ is exactly the inclusion
of finite-image functions. In particular $n\mapsto n$ is omitted.
No mapping-space theorem for unmarked vertical hammocks has been
used in any part of this proof.

The independent program \texttt{f\_fence\_paths.py} records the
complete four-arrow circle table and the covering-star checks,
all relative contractions of fences of radii one through eight,
and the full symmetric path stages $n=0,1,2$. Their object counts
are $1,29,985$, their arrow counts are $1,71,7573$, and all
$1,225,102471$ composable triples respectively satisfy the
category equations. Their connected components have exactly the
computed winding numbers. Sixteen explicit excursion maps are
also checked; the preceding formula proves the obstruction for
every height, independently of this finite range.
\end{proof}
